\section{QPU Experimental Protocol and Simulated Results}
\label{sec:qpu}

This section specifies the complete experimental protocol for execution on a Quandela Belenos (or equivalent) photonic QPU. The protocol is designed to be executed by Laurent and the Quandela team; simulated results are provided for reference. \textbf{All QPU results in this version are simulated; experimental data will be inserted upon QPU execution.}

\subsection{Hardware Requirements}

\begin{itemize}
\item \textbf{QPU:} Quandela Belenos (12-mode universal linear optical processor) or equivalent.
\item \textbf{Modes:} $n=4,6,8$ (start with $n=4$).
\item \textbf{Photons:} $k=n/2$ indistinguishable photons from a heralded single-photon source.
\item \textbf{Input state:} $|1^k 0^{n-k}\rangle$ (first $k$ modes occupied).
\item \textbf{Detection:} Photon-number-resolving (PNR) detectors on all $n$ modes.
\item \textbf{Shots:} 10,000--100,000 per circuit configuration.
\end{itemize}

\subsection{Calibration Phase (Mandatory)}

Before any optimization, the following calibrations must be performed and recorded:

\begin{enumerate}
\item \textbf{HOM calibration:} Prepare $|1,1,0,0\rangle$ on modes 0,1; apply 50:50 beamsplitter; measure $P_{\mathrm{CF}}$. Target: $P_{\mathrm{CF}} < 0.01$. Records indistinguishability and PNR calibration.
\item \textbf{Loss characterization:} Measure per-mode transmission $\eta_i$ using coherent state probes.
\item \textbf{Dark count rate:} Measure with vacuum input.
\item \textbf{Phase stability:} Measure phase drift over typical experiment duration.
\end{enumerate}

\subsection{Phase A: Haar Baseline}

For each $n \in \{4,6,8\}$:
\begin{enumerate}
\item Generate 100 Haar-random unitaries (seeded for reproducibility).
\item For each unitary:
  \begin{itemize}
  \item Compile to Clements mesh parameters.
  \item Upload to QPU.
  \item Prepare $|1^k 0^{n-k}\rangle$.
  \item Measure 1,000 shots.
  \item Record: $P_{\mathrm{CF}}$, $P_{\mathrm{success}}$, $P_{\mathrm{opt|CF}}$.
  \end{itemize}
\item Compute mean/std of $P_{\mathrm{CF}}$, $P_{\mathrm{success}}$.
\item Compare to theoretical Haar prediction $\mathbb{E}[P_{\mathrm{CF}}] = \binom{n}{k}/\binom{n+k-1}{k}$.
\end{enumerate}

\subsection{Phase B: Variational Optimization}

For each graph family $\times$ $n$:
\begin{enumerate}
\item \textbf{Graph instances:} Path, cycle, complete, bipartite, star, ER($p=0.5$), RR($d=3$), planted partition.
\item \textbf{Optimization:}
  \begin{itemize}
  \item Parametrize Clements mesh (depth $n$).
  \item Objective: Maximize $P_{\mathrm{success}}$ (exact distribution, no shots).
  \item Optimizer: COBYLA, 200 iterations (201 distribution evaluations).
  \item Seeds: 3 random restarts (seeds 42, 43, 44).
  \item Optional: Graph-informed initialization from spectral partitioning.
  \end{itemize}
\item \textbf{Validation:} Evaluate optimized circuit with 10,000 shots.
\item \textbf{Record:} $P_{\mathrm{CF}}$, $P_{\mathrm{success}}$, $P_{\mathrm{opt|CF}}$, expected shots, circuit parameters.
\end{enumerate}

\subsection{Phase C: Baseline Comparisons}

For each optimized circuit:
\begin{enumerate}
\item \textbf{Penalty baseline:} Same graph, $\lambda=10$, 200 random logit evaluations, 10,000 shots each.
\item \textbf{Distinguishable control:} Same optimized unitary, distinguishable photons (temporal delay $> \tau_c$ or spectral distinguishability).
\item \textbf{Uniform feasible:} Random balanced partitions (classical).
\end{enumerate}

\subsection{Phase D: Data Collection and Metrics}

For each experiment, record:
\begin{itemize}
\item Raw shot data (occupation numbers per shot) in CSV format.
\item Circuit parameters (mesh angles) in JSON.
\item Summary metrics with Wilson 95\% CI:
  \begin{itemize}
  \item $P_{\mathrm{CF}}$, $P_{\mathrm{success}}$, $P_{\mathrm{opt|CF}}$
  \item Expected shots $N_{\mathrm{shots}} = 1/P_{\mathrm{success}}$ with CI
  \item Conditional energy $\mathbb{E}[C|\mathrm{CF}]$
  \end{itemize}
\end{itemize}

\subsection{Simulated QPU Results (Reference)}

The following tables show simulated results for the protocol, assuming ideal hardware ($\eta=1$, $\xi=1$). These will be replaced with experimental data.

\begin{table}[ht]
\centering
\caption{Simulated QPU results for $n=4$ (10,000 shots per circuit).}
\label{tab:qpu_n4}
\begin{tabular}{lccccc}
\toprule
Graph & $P_{\mathrm{CF}}$ & $P_{\mathrm{success}}$ & $P_{\mathrm{opt|CF}}$ & $N_{\mathrm{shots}}$ & Speedup vs Penalty \\
\midrule
path & 1.000 & 1.000 & 1.000 & 1.0 & 6.0$\times$ \\
cycle & 1.000 & 0.999 & 0.999 & 1.0 & 3.5$\times$ \\
complete & 1.000 & 1.000 & 1.000 & 1.0 & 2.5$\times$ \\
bipartite & 1.000 & 1.000 & 1.000 & 1.0 & 4.4$\times$ \\
star & 1.000 & 1.000 & 1.000 & 1.0 & 2.5$\times$ \\
\bottomrule
\end{tabular}
\end{table}

\begin{table}[ht]
\centering
\caption{Simulated QPU results for $n=6$ (10,000 shots per circuit).}
\label{tab:qpu_n6}
\begin{tabular}{lccccc}
\toprule
Graph & $P_{\mathrm{CF}}$ & $P_{\mathrm{success}}$ & $P_{\mathrm{opt|CF}}$ & $N_{\mathrm{shots}}$ & Speedup vs Penalty \\
\midrule
path & 0.993 & 0.982 & 0.989 & 1.0 & 5.0$\times$ \\
cycle & 0.333 & 0.290 & 0.870 & 3.5 & 3.7$\times$ \\
complete & 0.992 & 0.992 & 1.000 & 1.0 & 2.9$\times$ \\
bipartite & 0.984 & 0.976 & 0.992 & 1.0 & 3.1$\times$ \\
star & 0.992 & 0.992 & 1.000 & 1.0 & 2.9$\times$ \\
ER & 1.000 & 0.987 & 0.987 & 1.0 & 32.9$\times$ \\
RR & 0.998 & 0.982 & 0.984 & 1.0 & 70.1$\times$ \\
planted & 0.994 & 0.988 & 0.995 & 1.0 & 29.1$\times$ \\
\bottomrule
\end{tabular}
\end{table}

\begin{table}[ht]
\centering
\caption{Simulated QPU results for $n=8$ (10,000 shots per circuit).}
\label{tab:qpu_n8}
\begin{tabular}{lccccc}
\toprule
Graph & $P_{\mathrm{CF}}$ & $P_{\mathrm{success}}$ & $P_{\mathrm{opt|CF}}$ & $N_{\mathrm{shots}}$ & Speedup vs Penalty \\
\midrule
path & 0.733 & 0.537 & 0.732 & 1.9 & 6.6$\times$ \\
cycle & 0.322 & 0.209 & 0.651 & 4.8 & 1.6$\times$ \\
complete & 0.852 & 0.852 & 1.000 & 1.2 & 3.0$\times$ \\
bipartite & 0.909 & 0.831 & 0.915 & 1.2 & 6.9$\times$ \\
star & 0.852 & 0.852 & 1.000 & 1.2 & 3.0$\times$ \\
ER & 0.087 & 0.060 & 0.687 & 16.8 & 15.0$\times$ \\
RR & 0.810 & 0.581 & 0.717 & 1.7 & 6.1$\times$ \\
planted & 0.733 & 0.537 & 0.732 & 1.9 & 6.6$\times$ \\
\bottomrule
\end{tabular}
\end{table}

\subsection{Resource Estimate}

\begin{table}[ht]
\centering
\caption{QPU resource estimate.}
\label{tab:qpu_resources}
\begin{tabular}{lcccc}
\toprule
Phase & Circuits & Shots/Circuit & Total Shots & Time (est.) \\
\midrule
Calibration & 1 & 10,000 & 10,000 & 1 hr \\
Haar baseline & 300 & 1,000 & 300,000 & 6 hrs \\
Variational & 24 & 10,000 & 240,000 & 8 hrs \\
Penalty baseline & 24 & 10,000 & 240,000 & 8 hrs \\
Distinguishable & 24 & 10,000 & 240,000 & 8 hrs \\
\midrule
Total & 597 & & 1,030,000 & $\sim$31 hrs \\
\bottomrule
\end{tabular}
\end{table}

\subsection{Success Criteria}

The QPU experiment is considered successful if:
\begin{enumerate}
\item HOM calibration yields $P_{\mathrm{CF}} < 0.01$.
\item Measured $P_{\mathrm{success}}$ for variational circuits exceeds Haar baseline by $>5\%$ (absolute) for $n=4,6$.
\item Measured $P_{\mathrm{success}}$ for variational circuits exceeds penalty baseline by $>2\times$ for $n=4,6$.
\item Bosonic advantage ($P_{\mathrm{success}}^{\mathrm{boson}} > P_{\mathrm{success}}^{\mathrm{dist}}$) observed for $n=8$ on at least one structured graph (RR or planted).
\end{enumerate}

If these criteria are not met, the results will be reported as a negative result with full data transparency.
